% Lean module: MonadFormal.Information.Testing.Exponent
\subsubsection{Optimal error exponent}
\label{sec:chernoff-exponent}

For $n\ge1$, define the rate
$r_n=-\log P_{\mathrm{err}}^*(n)/n$, with $r_n=+\infty$ when
the error is zero. Logarithms, and hence rates, are in nats.
\leandef{MonadFormal.Information.binaryErrorRate}

\begin{theorem}[Chernoff error exponent]
\label{thm:chernoff-exponent}
For two probability laws on a finite alphabet, independent identically
distributed observations and fixed strictly positive priors $a,b$,
\begin{equation}
 \lim_{n\to\infty}-\frac1n\log P_{\mathrm{err}}^*(n)
 =\ChernoffInfo(P,Q).
 \label{eq:error-exponent}
\end{equation}
The equality holds in the extended real line. Equal priors, identical
supports and an interior optimizing exponent are not required.
\end{theorem}
\begin{proof}
If $S=\varnothing$, \Cref{prop:chernoff-support} gives zero error
for every $n\ge1$ and $\ChernoffInfo(P,Q)=+\infty$.
Otherwise choose $s_*$ as above. Then $z>0$, and
\Cref{lem:chernoff-support-extension} gives
$\ChernoffInfo(P,Q)=-\log z$.
Put $c=\min\{a,b\}/2>0$. For each fixed $\varepsilon>0$,
the variance condition in \Cref{thm:chernoff-converse} holds for
all sufficiently large $n$. Its positive lower bound ensures
that taking ordinary real logarithms is legitimate, and yields
\[
 -\log z\le r_n\le-\log z+\varepsilon-\frac{\log c}{n}.
\]
The constant term $\log c/n$ tends to zero. For any neighborhood
of $-\log z$, first choose a smaller positive $\varepsilon$ and
then a sufficiently large $n$; the displayed inequalities place
$r_n$ in that neighborhood. This proves convergence and, by
embedding the real line, the claimed extended-valued statement.
\leanproof{MonadFormal.Information.rate_ge_of_error_le,MonadFormal.Information.rate_le_of_error_ge,MonadFormal.Information.tendsto_rate_of_exponential_bounds,MonadFormal.Information.binaryError_iid_eventually_ge,MonadFormal.Information.binaryError_iid_eventually_pos,MonadFormal.Information.chernoffInfo_eq_log_supportPartition,MonadFormal.Information.tendsto_binaryError_log_rate,MonadFormal.Information.binaryErrorRate_eq_coe,MonadFormal.Information.tendsto_binaryErrorRate}
\end{proof}

This is the classical binary-testing exponent
\cite{chernoff1952measure,cover2006elements}. It identifies the
exponential rate, not a finite-sample prefactor or an operational
channel-coding rate. Fixed positive priors disappear from the limit
because their logarithmic contribution is divided by $n$; priors
that vanish or depend on $n$ fall outside the theorem.
